\section{The auxiliary modes and the common kernel}
\label{app:darboux}
This appendix proves proposition~\ref{prop:modes}. We again work on the Hilbert
space $\HH=L^2((0,K),\dd s/K)$, with $0<q<1$.
We retain $\omega_r,\xi_r$ from section~\ref{sec:unflavoured}
and put $L=\pi K'/K$, $\dd\nu=\dd s/K$.

\subsection{The Fourier coefficients}
For $j\in\mathbb Z$, define
\begin{equation}
  \begin{aligned}
    \widehat c_j&=\int_0^K\cn(2s)\cos(j\pi s/K)\dd\nu(s),&
    \widehat s_j&=\int_0^K\sn(2s)\sin(j\pi s/K)\dd\nu(s),\\
    \widehat d_j&=\int_0^K\dn(2s)\cos(j\pi s/K)\dd\nu(s),&
    \widehat e_j&=\int_0^K k\sn(2s)\sin(j\pi s/K)\dd\nu(s),\\
    \widehat t_j&=\int_0^K\tan\varphi(s)\sin(j\pi s/K)\dd\nu(s).
  \end{aligned}
  \label{eq:fourier-coefficients}
\end{equation}
Put $\alpha=\pi/(2kK)$ and $\beta=\pi/(2K)$. Then
\begin{equation}
  \begin{array}{c|cc}
    &j\text{ odd}&j\text{ even}\\\hline
    \widehat c_j&\alpha\sech(jL/2)&0\\
    \widehat s_j&\alpha\csch(jL/2)&0\\
    \widehat d_j&0&\beta\sech(jL/2)\\
    \widehat e_j&\beta\csch(jL/2)&0\\
    \widehat t_j&2\alpha\csch(jL)&0
  \end{array}
  \label{eq:fourier-table}
\end{equation}
The even column includes $j=0$, in particular $\widehat d_0=\beta$. Negative
indices are included with the parity of $\sech$ and $\csch$.

The first four rows follow by substituting $2s$ into
\eqref{eq:sn-fourier}--\eqref{eq:dn-fourier} and integrating ordinary
trigonometric modes. To obtain the last row, note that
\begin{equation}
  \tan\varphi(s)=\frac{k\sn s\cn s}{\dn s}
  =-\frac1k\frac{\dd}{\dd s}\log\dn s.
\end{equation}
The theta product of $\dn s$, with $v=\pi s/K$, gives
\begin{equation}
  \log\dn s=\text{constant}
  +4\sum_{\substack{j\geq1\\j\text{ odd}}}
  \frac{q^j}{j(1-q^{2j})}\cos(jv).
\end{equation}
The series and its derivative converge uniformly on $[0,K]$. Differentiating
gives
\begin{equation}
  \tan\varphi(s)=\frac{2\pi}{kK}
  \sum_{\substack{j\geq1\\j\text{ odd}}}\csch(jL)\sin(j\pi s/K),
\end{equation}
which proves the fifth row.

\subsection{Orthonormality}
The duplication formulas applied to \eqref{eq:half-angles} imply
\begin{equation}
  \begin{aligned}
    \cos2\theta&=\cn(2s),&\sin2\theta&=\sn(2s),\\
    \cos2\varphi&=\dn(2s),&\sin2\varphi&=k\sn(2s).
  \end{aligned}
\end{equation}
Write $a_r=\cosh\xi_r$, $b_r=\sinh\xi_r$. In the overlap
$\langle\hat{\Phi}_r|\hat{\Phi}_\ell\rangle$, the constant part of
$\hat{\Phi}^*_r(s)\hat{\Phi}_\ell(s)$ integrates to
\begin{equation}
  \frac{a_ra_\ell+b_rb_\ell}
  {\sqrt{\cosh(2\xi_r)\cosh(2\xi_\ell)}}\delta_{r\ell}
  =\delta_{r\ell}.
\end{equation}
Set
\begin{equation}
  d=r-\ell,\qquad p=r+\ell+1,\qquad
  \delta=\xi_r-\xi_\ell=dL/2,\qquad
  \sigma=\xi_r+\xi_\ell=pL/2.
\end{equation}
The remaining part reduces to
\begin{equation}
  E^{\Phi}_{r\ell}=\frac1{\sqrt{\cosh(2\xi_r)\cosh(2\xi_\ell)}}\left[
    \widehat c_d\cosh\delta-\widehat s_d\sinh\delta
    -\widehat c_p\cosh\sigma+\widehat s_p\sinh\sigma\right].
\end{equation}
Since $d,p$ have opposite parity, either the $d$ terms vanish and the $p$ terms
cancel, or conversely. Thus $E^{\Phi}_{r\ell}=0$ by \eqref{eq:fourier-table}.
For $\langle\hat{\Psi}_r|\hat{\Psi}_\ell\rangle$ the corresponding correction is
\begin{equation}
  E^{\Psi}_{r\ell}=\frac1{\sqrt{\cosh(2\xi_r)\cosh(2\xi_\ell)}}\left[
    \widehat d_d\cosh\delta+\widehat e_d\sinh\delta
    -\widehat d_p\cosh\sigma-\widehat e_p\sinh\sigma\right].
\end{equation}
If $d$ is even, the first and last terms are $\beta$ and $-\beta$, with the
others zero. If $d$ is odd, the second and third cancel instead. This includes
$r=\ell$, using $\widehat d_0=\beta$. Consequently,
\begin{equation}
  \langle\hat{\Phi}_r|\hat{\Phi}_\ell\rangle
  =\langle\hat{\Psi}_r|\hat{\Psi}_\ell\rangle
  =\delta_{r\ell}.
  \label{eq:mode-orthogonality}
\end{equation}

\subsection{Completeness}
We first record the following fact needed for completeness.
\begin{lemma}\label{lem:phase}
  Let $\mu\in\{1,\tfrac12\}$ and let $\delta\in C([0,K],\mathbb R)$ satisfy
  $\delta(0)=\delta(K)=0$. The synthesis map
  \begin{equation}
    F_\delta a=\sqrt2\sum_{r\geq0}a_r
    \sin\!\left(\frac{(r+\mu)\pi s}{K}+\delta(s)\right)
  \end{equation}
  extends from finite sequences to a bounded Fredholm map $\ell^2(\mathbb
    N_0)\to\HH$ of index zero.
\end{lemma}
\begin{proof}
  Extend $\delta$ oddly to $(-K,K)$ and periodically across its endpoints. For
  $\mu=1$ use the integer lattice
  $\left\{\frac{j\pi}{K}:j\in\mathrm{Z}\right\}$; for $\mu=1/2$ use the
  half-integer lattice
  $\left\{\frac{\left(j+\frac{1}{2}\right)\pi}{K}:j\in\mathrm{Z}\right\}$.
  Let $P_+,P_-$ project onto its positive and negative frequencies,
  respectively. On the odd subspace, the operator
  \begin{equation}
    T_\delta=M_{e^{\ii\delta}}P_++M_{e^{-\ii\delta}}P_-
  \end{equation}
  preserves oddness and sends $\sin(\lambda s)$ to $\sin(\lambda s+\delta(s))$.
  Under restriction to $(0,K)$, it represents $F_\delta$. Consider
  \begin{equation}
    B_\delta=P_+M_{e^{-\ii\delta}}+P_-M_{e^{\ii\delta}}.
  \end{equation}
  Then the products $B_\delta T_\delta$ and $T_\delta B_\delta$ are the identity
  modulo compact operators. Indeed, the errors are cross-frequency terms and
  commutators of Fourier projections with continuous multiplication operators.

  For a trigonometric-polynomial such as $g_N(s)=\sum_{m=-N}^Nc_m\exp(im\pi
  s/K)$ the constituent operators in the error terms are finite rank. Since
  every continuous function $g$ with period $2K$ can be uniformly approximated
  by some $g_N$ the finite-rank property extends to compactness. The zero
  Fourier mode in the integer lattice contributes at most a finite-rank
  correction. Thus, from Atkinson's stability theorem the restriction to the odd
  subspace is Fredholm.

  Finally, replacing $\delta$ by $t\delta$, $0\leq t\leq1$, gives a
  norm-continuous Fredholm path from the ordinary sine synthesis operator, which
  is unitary. The index is therefore zero.
\end{proof}

For the normalised $\hat{\Phi}$ modes, expansion of the hyperbolic functions
gives
\begin{equation}
  \hat{\Phi}_r(s)=a_r^+\sin(\omega_rs+\theta(s))
  +a_r^-\sin(\omega_rs-\theta(s)),\qquad
  a_r^\pm=\frac{e^{\pm\xi_r}}{\sqrt{\cosh(2\xi_r)}}.
\end{equation}
Here $a_r^+=\sqrt2+O(e^{-4\xi_r})$ and $a_r^-=O(e^{-2\xi_r})$. Since
\begin{equation}
  \omega_rs+\theta(s)=\frac{(r+1)\pi s}{K}
  +\delta_\Phi(s),\qquad
  \delta_\Phi(s)=\theta(s)-\frac{\pi s}{2K},
\end{equation}
and $\delta_\Phi$ vanishes at both endpoints, lemma~\ref{lem:phase} applies with
$\mu=1$. The difference between synthesis by the $\hat{\Phi}_r$ and by the
phase-shifted sines is square-summable,
\begin{equation}
  h_r(s)=\hat{\Phi}_r(s)
  -\sqrt{2}\sin\left(\frac{(r+1)\pi s}{K}+\delta_\Phi (s)\right),
  \quad \sum_{r=0}^\infty ||h_r||_\mathscr{H}^2 < \infty.
\end{equation}
Therefore the difference is a Hilbert-Schmidt operator, and hence compact. By
the Fredholm stability theorem synthesis by $\hat{\Phi}_r$ is therefore Fredholm
and of index zero. By \eqref{eq:mode-orthogonality} it is also an isometry. Its
kernel and cokernel must both vanish, so it is surjective.

For $\hat{\Psi}_r(s)$ the same expansion has leading term
$\sqrt2\sin(\omega_rs-\varphi(s))$. The phase $-\varphi$ vanishes at both
endpoints, so lemma~\ref{lem:phase} applies with $\mu=1/2$. The same
square-summable perturbation argument proves that the
$\hat{\Psi}_r$ form a complete system.
Together this proves the basis assertion in proposition~\ref{prop:modes}.

\subsection{The cross identity}
Let $c_0=\sqrt{k/2}$ and $s_r(s)=\langle s|s_r\rangle=\sin(\omega_rs)$.
The weighted sine has wavefunction $\langle s|m_G s_r\rangle=m_G(s)s_r(s)$.
The definitions give four multiplier identities,
\begin{equation}
  \begin{aligned}
    m_C\cos\theta&=c_0(1+\cn(2s)),&
    m_C\sin\theta&=c_0\sn(2s),\\
    m_B\cos\varphi&=c_0,&m_B\sin\varphi&=c_0\tan\varphi.
  \end{aligned}
\end{equation}
The matrix-element identity~\eqref{eq:bounded-coefficients} reduces to
\begin{equation}
  \sqrt{\gamma_\ell}\langle\hat{\Phi}_r|m_C s_\ell\rangle
  =\sqrt{\gamma_r}\langle m_B s_r|\hat{\Psi}_\ell\rangle.
  \label{eq:cross}
\end{equation}
After removing the common factor $2\sqrt{kK/\pi}$, the constant part of
the difference in~\eqref{eq:cross} is
$c_0(a_r-a_\ell)\langle s_r|s_\ell\rangle=0$.
For $d=r-\ell$, $p=r+\ell+1$, the remaining terms are
\begin{equation}
  \begin{aligned}
    &\frac{\sqrt{\pi}}{c_0\sqrt{kK}}
    \left(\frac{\langle\hat{\Phi}_r|m_C s_\ell\rangle}{\sqrt{\gamma_r}}
    -\frac{\langle m_B s_r|\hat{\Psi}_\ell\rangle}{\sqrt{\gamma_\ell}}\right)
    =a_r(\widehat c_d-\widehat c_p)
    +b_r(\widehat s_p-\widehat s_d)
    +b_\ell(\widehat t_p+\widehat t_d).
  \end{aligned}
\end{equation}
If $d$ is odd, this is
\begin{equation}
  \alpha\left[a_r\sech\delta-b_r\csch\delta
    +2b_\ell\csch(2\delta)\right]=0,
  \qquad\delta=\xi_r-\xi_\ell,
\end{equation}
since $a_r\sinh\delta-b_r\cosh\delta=-\sinh\xi_\ell$.
If $d$ is even, $p$ is odd and the expression is
\begin{equation}
  \alpha\left[-a_r\sech\sigma+b_r\csch\sigma
    +2b_\ell\csch(2\sigma)\right]=0,
  \qquad\sigma=\xi_r+\xi_\ell,
\end{equation}
since $-a_r\sinh\sigma+b_r\cosh\sigma=-\sinh\xi_\ell$.
This proves \eqref{eq:cross}, including the diagonal case.

The geometric decay of $\sqrt{\gamma_r}$ and the uniform bounds on the
normalised modes and multipliers give uniform convergence of the
series~\eqref{eq:darboux-series} on $[0,K]^2$, completing the proof of
proposition~\ref{prop:modes}.
